\chapter{Smart Contracts and Crypto Engineering}\label{ch:iv3:smart-contracts-and-crypto-engineering}

\begin{ivtrains}
What a trading engineer at a crypto firm should know: how the EVM charges for work, the classic contract flaws and their fixes, oracles, MEV, and the plumbing of an exchange. Each vulnerable contract is shown with its flaw and its repair, on a local in-memory EVM or a Python model, and none refers to anything deployed. Solidity here is version 0.8.28, compiled by solc for the Cancun fork, unoptimised.
\end{ivtrains}

\begin{question}[$\star$ \iqroles{QD, Dev} \iqfirm{crypto} \iqlevel{junior}]\label{iq:iv3:smart-contracts-and-crypto-engineering:1}
Alice holds two unspent coins, worth 500 and 300 units, on a chain that uses unspent transaction outputs. She pays 600 with a fee of 10. In Python, write the spend: what goes in, what comes out, and what does she hold after? How would the same payment look on a chain with accounts? Name one thing each model makes easier.
\end{question}

\begin{question}[$\star$ \iqroles{QD, Dev, QT} \iqfirm{crypto} \iqlevel{junior}]\label{iq:iv3:smart-contracts-and-crypto-engineering:2}
On Ethereum, a plain transfer uses 21\,000 gas. The base fee is 20 gwei, and you add a priority fee of 1 gwei. One gwei is $10^{-9}$ ETH. What is the fee in ETH, and who gets each part? If the block is full, how far can the next base fee rise?
\end{question}

\begin{question}[$\star$ \iqroles{QD, Dev} \iqfirm{crypto} \iqlevel{junior}]\label{iq:iv3:smart-contracts-and-crypto-engineering:3}
In \cref{lst:iv3:cr:fixed}, in Solidity 0.8, what do \texttt{divFirst(1, 3, 3)}, \texttt{mulFirst(1, 3, 3)}, \texttt{wmul(1.5e18, 2e18)} and \texttt{wmul(1, 1)} return? What does \texttt{mulFirst} do on $(2^{255}, 2, 2)$, although the true answer fits in 256 bits? Why does Solidity have no floating point, and how do you write 1.5 instead?
\omcode{code/interviews-3/08-smart-contracts-and-crypto-engineering/solidity/Fixed.sol}{4}{18}{Integer arithmetic.\protect\hfil\protect\break}\label{lst:iv3:cr:fixed}
\end{question}

\begin{question}[$\star$ \iqroles{QD, Dev} \iqfirm{crypto} \iqlevel{junior}]\label{iq:iv3:smart-contracts-and-crypto-engineering:4}
A new deployment of \cref{lst:iv3:cr:locations} has \texttt{p.x} equal to 1. What does \texttt{viaMemory()} return, and \texttt{viaStorage()} afterwards? What do \texttt{memory}, \texttt{storage} and \texttt{calldata} mean, and how long does each one last? Which is the expensive one?
\omcode{code/interviews-3/08-smart-contracts-and-crypto-engineering/solidity/Locations.sol}{4}{19}{Data locations.\protect\hfil\protect\break}\label{lst:iv3:cr:locations}
\end{question}

\begin{question}[$\star$ \iqroles{QD, Dev} \iqfirm{crypto} \iqlevel{junior}]\label{iq:iv3:smart-contracts-and-crypto-engineering:5}
What is an account's nonce on Ethereum? What stops someone from broadcasting your signed transfer a second time? Your account's next nonce is 0, and your transactions with nonces 0, 1 and 3 are pending. Which can be included, and what happens to the last one? What happens if two transactions share a nonce?
\end{question}

\begin{question}[$\star$ \iqroles{QD, Dev, QT} \iqfirm{crypto} \iqlevel{junior}]\label{iq:iv3:smart-contracts-and-crypto-engineering:6}
Why does an exchange not credit a deposit as soon as it sees it in a block? On Ethereum proof of stake, with 12-second slots and 32 slots in an epoch, finality takes about two epochs. How long is that? What does reversing a finalised block cost an attacker, and what about a block that is only one slot deep?
\end{question}

\begin{question}[$\star\star$ \iqroles{QD, Dev} \iqfirm{crypto} \iqlevel{junior, experienced}]\label{iq:iv3:smart-contracts-and-crypto-engineering:7}
On \cref{lst:iv3:cr:slot}, under the Cancun fork's gas rules, the program sends four separate transactions: \texttt{set(1)}, \texttt{set(2)}, \texttt{set(2)} and \texttt{set(0)}, starting from \texttt{v} equal to 0. The call's own overhead is the same each time. Taking out that overhead, what does each write cost in gas, and is anything refunded? Why are the costs what they are?
\omcode{code/interviews-3/08-smart-contracts-and-crypto-engineering/solidity/Slot.sol}{4}{10}{One storage slot.\protect\hfil\protect\break}\label{lst:iv3:cr:slot}
\end{question}

\begin{question}[$\star\star$ \iqroles{QD, Dev} \iqfirm{crypto} \iqlevel{junior, experienced}]\label{iq:iv3:smart-contracts-and-crypto-engineering:8}
How many storage slots does each contract in \cref{lst:iv3:cr:packing} use? Draw slot 0 of \texttt{Tight}. What is the saving worth at deployment? When would you not reorder your variables? Why can you not simply reorder them in a new version of an upgradeable contract?
\omcode{code/interviews-3/08-smart-contracts-and-crypto-engineering/solidity/Packing.sol}{4}{14}{Two layouts.\protect\hfil\protect\break}\label{lst:iv3:cr:packing}
\end{question}

\begin{question}[$\star\star$ \iqroles{QD, Dev} \iqfirm{crypto} \iqlevel{junior, experienced}]\label{iq:iv3:smart-contracts-and-crypto-engineering:9}
What do \texttt{checkedDec(0)}, \texttt{uncheckedDec(0)} and \texttt{narrow(300)} do in \cref{lst:iv3:cr:overflow}? Which of them does the compiler protect, and which not? Why would anyone write \texttt{unchecked}? Give one case in which it is safe, and one in which it is not.
\omcode{code/interviews-3/08-smart-contracts-and-crypto-engineering/solidity/Overflow.sol}{4}{16}{Overflow and narrowing.\protect\hfil\protect\break}\label{lst:iv3:cr:overflow}
\end{question}

\begin{question}[$\star\star$ \iqroles{QD, Dev} \iqfirm{crypto} \iqlevel{junior, experienced}]\label{iq:iv3:smart-contracts-and-crypto-engineering:10}
The local \texttt{Vault} in \cref{lst:iv3:cr:vault} holds 10 ETH of Bob's. \texttt{Attacker.attack()} is called with 1 ETH. How much do the vault and the attacker hold afterwards? Explain the flaw step by step. Give two ways to fix it, and say which you would pick. Does Bob's withdrawal still work after each fix?
\omcode{code/interviews-3/08-smart-contracts-and-crypto-engineering/solidity/Vault.sol}{4}{33}{A vault and a caller.\protect\hfil\protect\break}\label{lst:iv3:cr:vault}
\end{question}

\begin{question}[$\star\star$ \iqroles{QD, Dev} \iqfirm{crypto} \iqlevel{junior, experienced}]\label{iq:iv3:smart-contracts-and-crypto-engineering:11}
\texttt{Payout} in \cref{lst:iv3:cr:pay} pays a contract that refuses all payments. After \texttt{payUnchecked} with 1 ETH, what do \texttt{paid}, the recipient's balance and \texttt{Payout}'s balance show? What does \texttt{payChecked} do? What other faults does a raw \texttt{call} invite? A contract pays 100 recipients in a loop: why is that a poor design?
\omcode{code/interviews-3/08-smart-contracts-and-crypto-engineering/solidity/Pay.sol}{4}{21}{Paying someone who refuses.\protect\hfil\protect\break}\label{lst:iv3:cr:pay}
\end{question}

\begin{question}[$\star\star$ \iqroles{QD, Dev} \iqfirm{crypto} \iqlevel{junior}]\label{iq:iv3:smart-contracts-and-crypto-engineering:12}
The owner of the local \texttt{Wallet} in \cref{lst:iv3:cr:origin}, holding 5 ETH, is persuaded to call \texttt{Trap.claimPrize}. What happens with \texttt{byOrigin} true, and with it false? What is the difference between \texttt{tx.origin} and \texttt{msg.sender}? Why does the check not protect the owner?
\omcode{code/interviews-3/08-smart-contracts-and-crypto-engineering/solidity/Origin.sol}{4}{28}{Who is calling?\protect\hfil\protect\break}\label{lst:iv3:cr:origin}
\end{question}

\begin{question}[$\star\star$ \iqroles{QD, Dev} \iqfirm{crypto} \iqlevel{junior, experienced}]\label{iq:iv3:smart-contracts-and-crypto-engineering:13}
In a standard token, \texttt{approve(spender, m)} overwrites the allowance. An owner has allowed 100 and sends a transaction that lowers it to 60. The spender watches pending transactions and can have its own executed first. What is the most it can take? Show it in Python. How would you repair it, in the contract and in the client?
\end{question}

\begin{question}[$\star\star$ \iqroles{QD, Dev, QT} \iqfirm{crypto} \iqlevel{junior, experienced}]\label{iq:iv3:smart-contracts-and-crypto-engineering:14}
A constant-product pool with a 0.3\% fee holds 1000 ETH and 2\,000\,000 dollars of a stablecoin, each with 18 decimals. Write \texttt{amountOut} in Solidity using integers. How many dollars does 1 ETH buy, and how does that split between fee and price impact? Why does the division round down, and whom does that favour? How do you check that the contract agrees with a Python model?
\end{question}

\begin{question}[$\star\star$ \iqroles{QD, Dev, Risk} \iqfirm{crypto} \iqlevel{junior, experienced}]\label{iq:iv3:smart-contracts-and-crypto-engineering:15}
A lender values collateral at the spot price of a thin constant-product pool, at 1000 tokens and 1000 dollars, with a 0.3\% fee taken from the input, and lends up to 50\% of that value. Someone who owns 100 tokens, worth 100 dollars, can borrow 9000 dollars within one transaction, at no cost, and puts all of it into the pool. In Python, model the sequence and find the profit. What is the profit if the lender uses a fair price instead? What makes the attack possible, and what are the defences?
\end{question}

\begin{question}[$\star\star$ \iqroles{QD, Dev} \iqfirm{crypto} \iqlevel{experienced}]\label{iq:iv3:smart-contracts-and-crypto-engineering:16}
Instead of reading a pool's spot price, a contract can keep a running sum of price times elapsed seconds. In Python, write that accumulator, and the average price between two readings. If the price spikes from 1 to 100 for one of 100 seconds, what is the average? What does an attacker now have to do? What does the window trade off? What does an arithmetic average do badly here?
\end{question}

\begin{question}[$\star\star$ \iqroles{QD, Dev} \iqfirm{crypto} \iqlevel{junior, experienced}]\label{iq:iv3:smart-contracts-and-crypto-engineering:17}
A local price feed in \cref{lst:iv3:cr:feed} last reported 200\,000 at time $T$. \texttt{careful} is called with a limit of one hour at $T + 60$ s and at $T + 7200$ s. Then the feed reports $-1$. What do \texttt{trusting} and \texttt{careful} return in each of these cases? What else should a careful consumer check? What should the protocol do when it refuses a price?
\omcode{code/interviews-3/08-smart-contracts-and-crypto-engineering/solidity/Feed.sol}{4}{25}{Two consumers of a price feed.\protect\hfil\protect\break}\label{lst:iv3:cr:feed}
\end{question}

\begin{question}[$\star\star$ \iqroles{QD, QT} \iqfirm{crypto} \iqlevel{junior, experienced}]\label{iq:iv3:smart-contracts-and-crypto-engineering:18}
A constant-product pool holds 1000 tokens and 2\,000\,000 dollars, with a 0.3\% fee taken from the input. A pending swap spends 20\,000 dollars, and it will revert if it receives more than 1\% less than quoted now. In Python, model a trader who buys first and sells straight after, and find the largest first trade that lets the victim through, and its profit. How does this vary with the limit? What is this called, and what can the user do about it? Why not set the limit at 0?
\end{question}

\begin{question}[$\star\star$ \iqroles{QD, Dev} \iqfirm{crypto} \iqlevel{experienced}]\label{iq:iv3:smart-contracts-and-crypto-engineering:19}
Under the 16-and-4 rule, a non-zero byte of transaction data costs 16 gas and a zero byte costs 4. In Python, compute the cost of a standard ABI-encoded \texttt{transfer(address, uint256)} of 1 token with 18 decimals to an address with no zero bytes, and of a tightly packed version: a 1-byte selector, the 20-byte address, and the amount with its leading zeros dropped. Why do rollups care so much? What makes up a fee on a layer 2, and what is different about the code you write for one?
\end{question}

\begin{question}[$\star\star$ \iqroles{QD, Dev} \iqfirm{crypto} \iqlevel{junior, experienced}]\label{iq:iv3:smart-contracts-and-crypto-engineering:20}
In Python, write a deposit processor for an exchange. It sees blocks with their parents, and it credits a deposit once it has $N$ confirmations on the best chain. A reorganisation can replace recent blocks, and the same transaction may reappear on the new branch. What makes a credit happen once only? What must happen if a reorganisation goes deeper than $N$? How do you test it?
\end{question}

\begin{question}[$\star\star$ \iqroles{QD, Dev} \iqfirm{crypto} \iqlevel{junior, experienced}]\label{iq:iv3:smart-contracts-and-crypto-engineering:21}
An exchange pays withdrawals from one hot account, using nonces in order. A transaction is stuck, as its fee is too low, and the operator wants to retry. In Python, write a sender and a model chain. Why must a retry keep its nonce, and not get a new one? What does a stuck nonce do to the requests behind it? How do you bump a fee? What must be stored before anything is broadcast?
\end{question}

\begin{question}[$\star\star\star$ \iqroles{QD, Dev} \iqfirm{crypto} \iqlevel{experienced}]\label{iq:iv3:smart-contracts-and-crypto-engineering:22}
In \cref{lst:iv3:cr:proxy}, someone calls \texttt{setX(0xdead)} through the proxy. Where does that write go? What does \texttt{impl()} return afterwards, and what does a call of \texttt{x()} through the proxy do then? Explain \texttt{delegatecall}. How do you lay out storage to avoid this? What must an upgraded implementation respect?
\omcode{code/interviews-3/08-smart-contracts-and-crypto-engineering/solidity/Proxy.sol}{4}{22}{A proxy and its implementation.\protect\hfil\protect\break}\label{lst:iv3:cr:proxy}
\end{question}

\begin{question}[$\star\star\star$ \iqroles{QD, Dev} \iqfirm{crypto} \iqlevel{experienced}]\label{iq:iv3:smart-contracts-and-crypto-engineering:23}
In \cref{lst:iv3:cr:readonly}, Bob has joined with 100 ETH and \texttt{Reader.run()} is called with another 100 ETH, on a local chain. What price does \texttt{Reader} see in \texttt{receive}, and what is it just before and after? No state of the pool is changed during the callback, so why is this a form of reentrancy? Why is a lock on \texttt{leave} alone not enough? What would a lender reading this price lose? Fix it.
\omcode{code/interviews-3/08-smart-contracts-and-crypto-engineering/solidity/ReadOnly.sol}{4}{36}{A price read in the middle of a call.\protect\hfil\protect\break}\label{lst:iv3:cr:readonly}
\end{question}

\begin{question}[$\star\star\star$ \iqroles{QD, Dev} \iqfirm{crypto} \iqlevel{experienced}]\label{iq:iv3:smart-contracts-and-crypto-engineering:24}
\Cref{lst:iv3:cr:relay} lets anyone submit a user's signed payment and pay the gas for it. The signer has funded 10 ETH in each deployment, and a signature pays Mallory 1 ETH. If Mallory submits it twice, and then once to another deployment of the same contract, what happens? What is missing from the signed message, and why does each missing field matter? Write the repair. Why can a user pay without holding gas this way?
\omcode{code/interviews-3/08-smart-contracts-and-crypto-engineering/solidity/Relay.sol}{4}{18}{A relayed payment.\protect\hfil\protect\break}\label{lst:iv3:cr:relay}
\end{question}

\begin{question}[$\star\star\star$ \iqroles{QD, Dev, Risk} \iqfirm{crypto} \iqlevel{experienced}]\label{iq:iv3:smart-contracts-and-crypto-engineering:25}
A key that controls funds must survive the loss of one holder, and resist the theft of one. In Python, write Shamir secret sharing, $t$ of $n$ over a prime field. Why do $t - 1$ shares say nothing? If each of 3 key holders is independently compromised with probability 1\%, what is the chance of a theft under 2-of-3, and what is it with one key? How does this differ from a multisignature contract and from threshold signing?
\end{question}

\begin{question}[$\star\star\star$ \iqroles{QD, Dev, Risk} \iqfirm{crypto} \iqlevel{experienced}]\label{iq:iv3:smart-contracts-and-crypto-engineering:26}
An exchange publishes the root of a Merkle tree of account balances, so that a user can check their own. In Python, write the tree, an inclusion proof and its check. Why do leaves and inner nodes get different prefixes? How long, at most, is a proof for 23 accounts? What does this prove, and what does it not? What else must be shown for a claim of solvency?
\end{question}

\begin{question}[$\star\star\star$ \iqroles{QD, QT, Risk} \iqfirm{crypto} \iqlevel{experienced}]\label{iq:iv3:smart-contracts-and-crypto-engineering:27}
A long perpetual position of 1 token at 2000, with isolated margin of 200 and a maintenance rate of 0.5\%, is closed when equity falls to maintenance times notional at the mark. In Python, find the liquidation price. Why use a mark price and not the last trade? The price gaps and the close fills at 1790: who pays? What happens if the insurance fund runs out?
\end{question}

\begin{question}[$\star\star\star$ \iqroles{QD, QT} \iqfirm{crypto} \iqlevel{experienced}]\label{iq:iv3:smart-contracts-and-crypto-engineering:28}
A constant-product pool holds 1000 tokens and 2\,000\,000 dollars, with a 0.3\% fee taken from the input, and a centralised venue quotes the token at 2020. Ignoring that venue's costs, derive the best size to buy from the pool and sell there, and find the profit. Below what outside price is there no trade at all? Write it in Python. What risks does the formula leave out, in execution and in inventory?
\end{question}

\begin{question}[$\star\star\star\star$ \iqroles{QD, Dev, LLE} \iqfirm{crypto} \iqlevel{experienced}]\label{iq:iv3:smart-contracts-and-crypto-engineering:29}
Design a centralised crypto exchange with deposits, withdrawals, limit orders and a live book. How are balances kept and held for open orders? In Python, write the balance core. Where does the matching engine meet the ledger? How do deposits, withdrawals and custody work, and how much goes in a hot wallet? What must never be possible, and how would you check that it never happened?
\end{question}

\begin{question}[$\star\star\star\star$ \iqroles{QD, QT, LLE} \iqfirm{crypto} \iqlevel{experienced}]\label{iq:iv3:smart-contracts-and-crypto-engineering:30}
Design a system that makes markets on an on-chain pool and hedges on centralised venues. What does it read from the chain, and at what stage of a block? How do reorganisations, failed transactions and gas change the design? How do you handle the time between a decision and inclusion? What is the state of a transaction you have sent, and how is your position kept? What would you monitor? What must a test environment be able to do?
\end{question}
